\vspace{0.5cm}

\textbf{Security and Identity Management}

The application's security is built upon \textbf{Azure Active Directory (Azure AD)}. Azure AD is an identity management system that acts as a central \textbf{Identity Provider (IdP)}. It not only handles user authentication but also serves as the authority for token issuance. The architecture implements a modern OAuth 2.0 and OpenID Connect flow. OAuth 2.0 is used for authorization, allowing applications to access resources on behalf of a user without knowing their password, while OpenID Connect is used for authentication and provides information about the user's identity.

The client side handles user interaction for signing in to the application and uses the Microsoft Authentication Library (MSAL) for orchestrating login popups and redirects against Azure AD. After the user successfully logs in, MSAL obtains a JWT access token and stores it securely in session storage. The diagram illustrates a two-way communication between the \textbf{browser} and \textbf{Azure AD} during this process.

All further communication with the backend, including \textbf{REST API} calls and real-time \textbf{SignalR} connections, is secured using this JWT token. On the client side, the token is automatically added to the Authorization header as a Bearer token to communicate with the backend API. The backend API, hosted on \textbf{Azure App Service}, is protected using Microsoft.Identity.Web library. This middleware intercepts each incoming request, validates the JWT signature, claims (such as issuer and audience) against Azure AD's OpenID Connect metadata endpoints, and establishes the user identity.

This validation step is shown in the diagram as a one-way communication from the \textbf{App Service} to \textbf{Azure AD}. After successful validation, controllers and SignalR hubs enforce access control using the \texttt{[Authorize]} attribute. The application also maps the token's object identifier claim to a specific user, allowing it to perform authorized actions.
